\chapter{ชุดโจทย์และข้อสอบประยุกต์ฟิสิกส์สำหรับวิทยาศาสตร์ชีวภาพ}
\label{app:bio_problems}
\index{ชุดโจทย์ประยุกต์ชีวภาพ}

\section{ชุดโจทย์คัดสรรข้อสอบกลางภาคและปลายภาค}

\begin{enumerate}
    \item \textbf{กลศาสตร์ของไหลและการดูดซึม:} ท่อลำเลียงไซเลมของต้นยางพารามีรัศมีเฉลี่ย $r = 25\,\mu\text{m}$ หากน้ำยางมีแรงตึงผิว $\gamma = 0.072\,\text{N/m}$ และมุมสัมผัส $\theta = 0^\circ$ จงคำนวณหาความสูงสูงสุดที่ของเหลวสามารถขึ้นไปได้ด้วยปรากฏการณ์คาพิลลารี (Capillary Action) กำหนดความหนาแน่น $\rho = 1,000\,\text{kg/m}^3$ และ $g = 9.8\,\text{m/s}^2$
    
    \item \textbf{เครื่องหมุนเหวี่ยงแยกตะกอน:} เครื่อง Centrifuge หมุนด้วยความถี่ $6,000\,\text{rpm}$ รัศมีการหมุนของหลอดทดลอง $r = 12\,\text{cm}$ จงคำนวณหาความเร่งสู่ศูนย์กลางในหน่วยเท่าของแรงโน้มถ่วงโลก ($RCF \times g$)
    
    \item \textbf{ไฟฟ้าชีวภาพ:} เยื่อหุ้มเซลล์ประสาทมีความหนา $d = 8.0\,\text{nm}$ มีความต่างศักย์ขณะพัก $V = -70\,\text{mV}$ จงคำนวณขนาดของสนามไฟฟ้าที่ตั้งฉากกับเยื่อหุ้มเซลล์ในหน่วยโวลต์ต่อเมตร ($\text{V/m}$) พร้อมอธิบายว่าสนามไฟฟ้านี้ส่งผลต่อการเคลื่อนที่ของไอโซเลท $\text{Na}^+$ และ $\text{K}^+$ อย่างไร
    
    \item \textbf{อุณหพลศาสตร์และการสังเคราะห์แสง:} ใบพืชได้รับพลังงานแสงอาทิตย์เฉลี่ย $800\,\text{W/m}^2$ หากใบพืชมีพื้นที่ผิวรวม $0.05\,\text{m}^2$ และมีประสิทธิภาพในการเปลี่ยนพลังงานแสงเป็นพลังงานเคมี $3.5\%$ จงคำนวณพลังงานเคมีที่สังเคราะห์ได้ในเวลา 1 ชั่วโมง
    
    \item \textbf{ทัศนศาสตร์กล้องจุลทรรศน์:} เลนส์ใกล้วัตถุของกล้องจุลทรรศน์มีค่า Numerical Aperture ($NA = 1.25$) เมื่อใช้แสงสีน้ำเงินความยาวคลื่น $\lambda = 450\,\text{nm}$ จงคำนวณขีดจำกัดในการแยกชัดของกล้อง (Resolving Limit) ตามเกณฑ์ของเรย์ลีห์
\end{enumerate}
